%% From https://www.latextemplates.com/template/medium-length-professional-cv

%%%%%%%%%%%%%%%%%%%%%%%%%%%%%%%%%%%%%%%%%
% Medium Length Professional CV
% LaTeX Template
% Version 3.0 (December 17, 2022)
%
% This template originates from:
% https://www.LaTeXTemplates.com
%
% Author:
% Vel (vel@latextemplates.com)
%
% Original author:
% Trey Hunner (http://www.treyhunner.com/)
%
% License:
% CC BY-NC-SA 4.0 (https://creativecommons.org/licenses/by-nc-sa/4.0/)
%
%%%%%%%%%%%%%%%%%%%%%%%%%%%%%%%%%%%%%%%%%

%----------------------------------------------------------------------------------------
%	PACKAGES AND OTHER DOCUMENT CONFIGURATIONS
%----------------------------------------------------------------------------------------

\documentclass[
	%a4paper, % Uncomment for A4 paper size (default is US letter)
	11pt, % Default font size, can use 10pt, 11pt or 12pt
]{resume} % Use the resume class

\usepackage[backend=biber,style=numeric]{biblatex}
\addbibresource{references.bib}

% \usepackage{ebgaramond} % Use the EB Garamond font
% \usepackage{fontspec}
% \newfontfamily\vfont{Caladea}
% \setmainfont{Caladea}

%------------------------------------------------

\name{Eric  (Chen-Chien) Chang} % Your name to appear at the top

% You can use the \address command up to 3 times for 3 different addresses or pieces of contact information
% Any new lines (\\) you use in the \address commands will be converted to symbols, so each address will appear as a single line.

% \address{123 Broadway \\ City, State 12345} % Main address

% \address{123 Pleasant Lane \\ City, State 12345} % A secondary address (optional)

% \address {alxtz.tw@gmail.com \\ github.com/alxtz \\ linkedin.com/in/alxtz}

\address{Email: {\color{teal}{\underline{\href{mailto:e850506@gmail.com}{e850506@gmail.com}}}} \\ Phone: {\color{teal}{\underline{+886-932786506}}}}
\address{Blog: {\color{teal}{\underline{\href{https://blog.unknowntpo.me/}{https://blog.unknowntpo.me/}}}}}
\address{Github: {\color{teal}{\underline{\href{https://github.com/unknowntpo/}{https://github.com/unknowntpo/}}}}}

% \item Test \href{https://www.example.com}{\color{teal}\underline{Visit Example}}
%----------------------------------------------------------------------------------------

\begin{document}

%----------------------------------------------------------------------------------------
%	DESCRIPTION SECTION
%----------------------------------------------------------------------------------------

\begin{rSection}{Summary}
\begin{rSubsection}{}{}{}{}
            \item Backend / Data Infrastructure Engineer specializing in \textbf{performance tuning}, \textbf{distributed systems}, and \textbf{open-source data platforms}.
            \item Core stack: \textbf{Go}, \textbf{Java}, \textbf{Rust}, \textbf{TypeScript}, \textbf{PostgreSQL}, \textbf{Kafka}, \textbf{Docker}, \textbf{Kubernetes}. Delivered measurable wins including first-byte latency \textbf{100s to 5s}, WebUI load time \textbf{60s to 1s}, and simulation jobs \textbf{2hr to 1min}.
	\end{rSubsection}
\end{rSection}

%----------------------------------------------------------------------------------------
%	OPEN SOURCE CONTRIBUTIONS SECTION
%----------------------------------------------------------------------------------------

\begin{rSection}{Open Source Contributions}
\begin{rSubsection}{Apache Kafka}{2024 - Present}{Apache Kafka Contributor. Tech Stack: {\color{teal}\textbf{Java, JVM}}}{{\color{teal}\underline{\href{https://github.com/apache/kafka}{GitHub}}}}
		\item Collaborating on Kafka Streams \textbf{KIP-770} input-buffer work ({\color{teal}\underline{\href{https://github.com/apache/kafka/pull/22458}{PR \#22458}}}) to add \path{input.buffer.max.bytes}, bounding buffering by bytes instead of record count for safer memory control under variable record sizes and partition assignments.
		\item Fixed \textbf{JVM container awareness} in producer performance tests and offset commit-before-partition-switch correctness ({\color{teal}\underline{\href{https://github.com/apache/kafka/pull/21584}{PR \#21584}}}).
	\end{rSubsection}

\begin{rSubsection}{Apache Gravitino}{2024 - Present}{Apache Gravitino Committer. Tech Stack: {\color{teal}\textbf{Java, Python, Docker, Kubernetes}}}{{\color{teal}\underline{\href{https://github.com/apache/gravitino}{GitHub}}}}
		\item Added PostgreSQL/MySQL integration-test container management and reproducible Docker image workflows; reduced Hive image size by \textbf{420MB} and ported Python client partition APIs.
	\end{rSubsection}

\begin{rSubsection}{Apache DataFusion and DataFusion Comet}{2025 - Present}{Contributor. Tech Stack: {\color{teal}\textbf{Rust, Scala, Apache Spark}}}{{\color{teal}\underline{\href{https://github.com/apache/datafusion}{GitHub}}}}
		\item Implemented Spark-compatible \texttt{str\_to\_map} in Rust ({\color{teal}\underline{\href{https://github.com/apache/datafusion/pull/20120}{PR \#20120}}}) and wired Comet's native SparkStrToMap UDF with SQL file coverage ({\color{teal}\underline{\href{https://github.com/apache/datafusion-comet/pull/3654}{PR \#3654}}}).
	\end{rSubsection}

\begin{rSubsection}{redis/rueidis}{Oct 2023 - Present}{Contributor. Tech Stack: {\color{teal}\textbf{Go, Redis}}}{{\color{teal}\underline{\href{https://github.com/redis/rueidis/pull/395}{PR \#395}}}}
		\item Implemented the go-redis API Adapter for teams migrating from go-redis.
	\end{rSubsection}
\end{rSection}

%----------------------------------------------------------------------------------------
%	SELECTED PROJECTS SECTION
%----------------------------------------------------------------------------------------

\begin{rSection}{Selected Projects}
\begin{rSubsection}{TWFoundry}{2026}{Real-time Taiwan transit data platform (hybrid cloud / homelab). Tech Stack: {\color{teal}\textbf{Kafka, Flink, Iceberg, ClickHouse, Airflow, Kubernetes (k0s), Cloudflare R2 / Workers / Pages, Docker, TypeScript}}}{{\color{teal}\underline{\href{https://github.com/unknowntpo/twfoundry}{GitHub}}}}
		\item Built a real-time transit data platform ingesting Taiwan TDX bus telemetry on a 5-minute cadence through \textbf{Kafka} into a partitioned \textbf{Cloudflare R2} data lake, running 24/7 on a self-hosted single-node \textbf{k0s Kubernetes} cluster.
		\item Architected a \textbf{Lambda} pipeline: a \textbf{Flink} speed layer for live anomaly detection (service-gap / bus-bunching) over the stream, plus an \textbf{Airflow}-orchestrated batch layer landing an \textbf{Apache Iceberg} lakehouse on R2 with \textbf{ClickHouse} OLAP modeling observations across service date, time slot, route, direction, vehicle, freshness, route progress, and stop context for headway / bunching and data-quality analysis.
		\item Served replayable map / timeline projection APIs from partitioned R2 artifacts via \textbf{Cloudflare Workers / Pages} Functions at the edge, decoupling public reads from the homelab cluster (no fallback on cache misses).
	\end{rSubsection}
\end{rSection}

%----------------------------------------------------------------------------------------
%	WORK EXPERIENCE SECTION
%----------------------------------------------------------------------------------------

\begin{rSection}{Experience}
    \begin{rSubsection}{Lawsnote}{January 2024 - December 2024}{Backend Software Engineer. Tech Stack: {\color{teal}\textbf{Typescript, NodeJS, Docker, MongoDB, Prometheus, Grafana}}}{Taipei, Taiwan}
        \item Implemented \textbf{HTTP streaming} for CPU-intensive computation results, reducing first-byte latency from \textbf{100s} to \textbf{5s}.
        \item Built a \textbf{MinIO events} and \textbf{Elasticsearch} monitoring path to detect missing data proactively and improve pipeline reliability.
        \item Profiled and fixed a \textbf{NAPI-rs} performance issue with \textbf{Pyroscope}, improving Node.js native package execution.
        \item Built an \textbf{LLM service} for extracting information from ambiguous regulatory text and automated government act crawling, re-crawl, and re-parse workflows.
    \end{rSubsection}

	\begin{rSubsection}{Mediatek (Contract)}{February 2022 - March 2023}{Backend Software Engineer. Tech Stack: {\color{teal}\textbf{Go, Linux, Docker, PostgreSQL}}}{Hsinchu, Taiwan}
            \item Maintained a low-power simulation and benchmarking data warehouse, improving WebUI page load time from \textbf{60s} to \textbf{1s} with \textbf{PostgreSQL index-only scans}: {\color{teal}{\underline{\href{https://blog.unknowntpo.me/idx-only-scan/}{Blog Post}}}}.
            \item Used \textbf{Pyroscope eBPF} to identify import bottlenecks, then applied \textbf{PostgreSQL bulk import} to reduce simulation job duration from \textbf{2hr} to \textbf{1min}.
            \item Designed a message-queue-based microservice architecture and operated a \textbf{5-node} service cluster for horizontally scalable long-running jobs.
            \item Reduced Go service memory consumption by \textbf{50\%} using \textbf{sync.Pool}: {\color{teal}{\underline{\href{https://blog.unknowntpo.me/syncpool/}{Blog Post}}}}.
            \item Built integration-test fixture factories and maintained observability with \textbf{Prometheus}, \textbf{Grafana}, \textbf{OpenTelemetry}, and \textbf{Pyroscope}.
	\end{rSubsection}

	\begin{rSubsection}{National Chen-Kung University}{September 2017 - July 2019}{Research Assistant. Tech Stack: \color{teal}\textbf{Python, C}}{Tainan, Taiwan}
            \item Conducted FDM 3D printing research across firmware, filament experiments, and hardware/software integration.
            \item Built experimental tooling with \textbf{Python}, \textbf{C}, \textbf{MicroPython}, and \textbf{Arduino}.
	\end{rSubsection}
\end{rSection}

%----------------------------------------------------------------------------------------
%	EDUCATION SECTION
%----------------------------------------------------------------------------------------

\begin{rSection}{Education}
National Taiwan Ocean University, Electrical Engineering coursework (Drop Out) \hfill 2013 - 2017
\end{rSection}

%----------------------------------------------------------------------------------------
%	SKILLS SECTION
%----------------------------------------------------------------------------------------

\begin{rSection}{Skills}
\textbf{Languages:} Go, Java, Rust, TypeScript, Python, C \\
\textbf{Data Warehouse / Lakehouse:} PostgreSQL, ClickHouse, Hive \\
\textbf{Streaming / Processing:} Kafka, Kafka Streams, Flink, Spark, SQL \\
\textbf{Runtime / Orchestration:} Airflow, Docker, Kubernetes, Linux \\
\textbf{Observability / Data Quality:} Prometheus, Grafana, OpenTelemetry, Pyroscope, Great Expectations \\
\textbf{DevOps / Applied AI:} Git, GitHub Actions, LLM service integration
\end{rSection}


%----------------------------------------------------------------------------------------
%	EXAMPLE SECTION
%----------------------------------------------------------------------------------------

%\begin{rSection}{Section Name}

	%Section content\ldots

%\end{rSection}

%----------------------------------------------------------------------------------------
\end{document}
